% Astra-2 / NS-46: complete residual, cutoff loss, exact parity and rank.
% Companion finite-arc Gram estimate is owned by the coordinator.
\subsection*{Dense finite radical blocks: complete residual and rank accounting}
\label{ns46-a2-sec-block}

\paragraph{Scope and input.}
All operators below are the complete physical $\mathsf W_{e^a}$,
with the domain, both pole terms, all prime rows and fixed smooth
cutoff $\chi_a$ of Proposition~\ref{prop:v135-growing-radical}.
Write $\phi=h/\|h\|_2$ for the even Gaussian arithmetic radical.
Choose fixed $r,c_0>0$ such that $|\widehat\phi(\xi)|^2\ge c_0$
for $|\xi|\le r$; these exist because $\widehat\phi(0)>0$.
Lemma~\ref{lem:ns46-finite-arc} supplies the following
uncut Gram lower bound for $m\ge2$ equally spaced translates,
spacing $d>0$, whenever $rd\le\pi/2$:
\begin{equation}\label{ns46-a2-eq-arc-input}
 K_m\succeq\gamma_m I,\qquad
 \gamma_m=\frac{c_0r}{m}
       \left(\frac{rd}{4\pi e}\right)^{2(m-1)}.
\end{equation}
The results here use this proved explicit bound as their input.
No lower bound for an infinite lattice with shrinking spacing is
being assumed. This is an analytic construction, not an interval
computation or a new finite-window certificate.

\begin{lemma}[Uniform complete tails for a fixed range of centers]
\label{ns46-a2-lem-fixed-centers}
Fix $T>0$. For $|t|\le T$, let
$g_{a,t}=\chi_a\phi(\cdot-t)$ and
$r_{a,t}=(1-\chi_a)\phi(\cdot-t)$. There is a constant $C_T$,
independent of $a,t$, such that, for all sufficiently large $a$,
\begin{align}
 \|r_{a,t}\|_{H^1}
  +\sup_x e^{2|x|}(|r_{a,t}(x)|+|r'_{a,t}(x)|)
 &\le C_Te^{-c_Te^{2a}},\notag\\
 \int_{|x|\ge a-1}|\phi(x-t)|^2\,dx
 &\le C_T^2e^{-2c_Te^{2a}},
 \quad c_T=\frac\pi4e^{-2(T+1)},
 \label{ns46-a2-eq-tails}\\
 \|\mathsf W_{e^a}g_{a,t}\|_2
 &\le C_T\sqrt a\,e^{2a}e^{-c_Te^{2a}}.
 \label{ns46-a2-eq-column-residual}
\end{align}
\end{lemma}
\begin{proof}
On the tail support $|x|\ge a-1$, one has
$|x-t|\ge a-1-T$. Put $X=e^{2|x-t|}$, so
$X\ge e^{2a-2(T+1)}$ and $e^{2|x|}\le e^{2T}X$.
The derivative bounds \eqref{eq:v135-gaussian-tail} and the uniformly
bounded derivatives of the fixed cutoff bound the weighted pointwise
tails by $C_TXe^{-\pi X/2}\le C_Te^{-\pi X/4}$.
Integrating the squared Gaussian bound after the substitution
$X=e^{2|x-t|}$ proves the $H^1$ and full tail-mass estimates, with
the same weaker exponent. Constants are uniform for all $|t|\le T$.

Polarized radicality gives, distributionally on $I_a$,
$\mathsf W_{e^a}g_{a,t}=-\mathcal W r_{a,t}$.
The archimedean multiplier is bounded by $C(1+|\xi|)$, so its
$L^2$ contribution is controlled by the $H^1$ tail. For $x\in I_a$,
each of the two full prime sums is bounded by
\[
 C_Te^{-c_Te^{2a}}e^{2|x|}
         \sum_{n\ge2}\frac{\Lambda(n)}{n^{5/2}}.
\]
This includes prime rows beyond $e^{2a}$; there is no finite-prime
truncation of the residual. Both pole integrals are bounded by
$C_Te^{-c_Te^{2a}}e^{|x|/2}$. Integration on $I_a$ gives
\eqref{ns46-a2-eq-column-residual}. Each $g_{a,t}$ is smooth and
compactly supported inside $I_a$, hence in the complete operator domain.
\end{proof}

\begin{proposition}[Gram loss and residual amplification]
\label{ns46-a2-prop-transfer}
Take any $m$ distinct centers in $[-T,T]$ whose uncut synthesis
Gram satisfies $K\succeq\gamma I$, $\gamma>0$.
Let $G$ synthesize the cutoff columns, $\Gamma=G^*G$, and
$\tau_a=C_Te^{-c_Te^{2a}}$. If
$m\tau_a^2\le\gamma/2$, then
\begin{equation}\label{ns46-a2-eq-transfer}
 \Gamma\succeq\gamma I/2,\qquad
 \operatorname{rank}G=m,\qquad
 \|\mathsf W_{e^a}P_G\|
 \le \sqrt{2m/\gamma}\,
         C_T\sqrt a\,e^{2a}e^{-c_Te^{2a}}.
\end{equation}
Here $P_G=G\Gamma^{-1}G^*$ is the ordinary orthogonal projection.
\end{proposition}
\begin{proof}
The Gram loss is positive semidefinite and
\[
 \operatorname{tr}(K-\Gamma)
 =\sum_j\int(1-\chi_a^2)|\phi(x-t_j)|^2\,dx
 \le m\tau_a^2.
\]
The full tail mass in \eqref{ns46-a2-eq-tails} is essential here:
one cannot replace $1-\chi_a^2$ by $(1-\chi_a)^2$ as an upper bound.
Thus $\Gamma\succeq\gamma I/2$. Since $G\Gamma^{-1/2}$ is an
isometry onto its range, the complete column residuals give
\[
 \|\mathsf W_{e^a}P_G\|
 =\|\mathsf W_{e^a}G\Gamma^{-1/2}\|
 \le\sqrt{2/\gamma}
       \left(\sum_j\|\mathsf W_{e^a}g_{a,t_j}\|^2\right)^{1/2}.
\]
This proves the displayed bound without a midpoint inverse or a
quadratic-value substitute for the operator residual.
\end{proof}

\begin{corollary}[A block of order $\lambda^2/\log\lambda$]
\label{ns46-a2-cor-large-block}
Put $\lambda=e^a$ and, for sufficiently large $a$, choose
\[
 J=\left\lfloor\frac{\lambda^2}{40000a}\right\rfloor,
 \quad m=2J+1,\quad d=1/J,\quad t_j=j/J\quad(-J\le j\le J).
\]
Using \eqref{ns46-a2-eq-arc-input}, the cutoff synthesis has
rank $m$ and its ordinary projection satisfies
\begin{equation}\label{ns46-a2-eq-large-residual}
 \|\mathsf W_{e^a}P_G\|\le\eta_a,
 \qquad \eta_a=C\exp(-\lambda^2/2000).
\end{equation}
The even and odd parts have exact dimensions $J+1$ and $J$,
and inherit the same residual bound. All constants and the eventual
threshold are independent of $a$; no numerical threshold is certified.
\end{corollary}
\begin{proof}
Use $T=1$. Then $c_T=\pi/(4e^4)>1/1000$.
Absorbing the polynomial prefactor in
\eqref{ns46-a2-eq-column-residual} gives a column residual at most
$Ce^{-\lambda^2/1000}$, and the full tail norm has the same bound.
For large $a$, $rd\le\pi/2$, $m\le\lambda^2/(10000a)$ and
$\log(4\pi e/(rd))\le3a$. Hence
\[
 \gamma_m^{-1/2}\le C\sqrt m\,
       \exp(3ma)\le C\sqrt m\exp(3\lambda^2/10000).
\]
In particular $m\tau_a^2/\gamma_m\to0$, so
Proposition~\ref{ns46-a2-prop-transfer} applies. Its resulting
bound is at most
$Cm\exp(-7\lambda^2/10000)\le C\exp(-\lambda^2/2000)$.

Reflection sends the $j$th cutoff column to the $-j$th column.
The coefficient embeddings consisting of the central column and
$(e_j+e_{-j})/\sqrt2$, or of $(e_j-e_{-j})/\sqrt2$, are
isometries of dimensions $J+1$ and $J$. The full Gram lower bound
makes both restricted syntheses injective. Their images are exactly
the two parity parts and their projections lie under $P_G$.
\end{proof}

\paragraph{Exact source accounting.}
An even source imposes no condition on the $J$ odd columns. In the
even block its orthogonality condition removes at most one dimension:
the resulting dimension is $J$ if its projection into that block is
nonzero, and $J+1$ otherwise. The lower dimension $J$ is enough for
the conclusion below; no identification of the cutoff column with
the actual source or unproved nonzero-overlap claim is needed.

\begin{corollary}[Necessary removal rank for positive coercivity]
\label{ns46-a2-cor-rank}
For every physical subspace $F_a$ with $\dim F_a<m$, there is a
unit $v_a\in\operatorname{Ran}G\cap F_a^\perp$ with
\[
 \|\mathsf W_{e^a}v_a\|\le\eta_a,
 \qquad |q_a[v_a]|\le\eta_a.
\]
Thus every positive coercivity constant on $F_a^\perp$ is at most
$\eta_a$. In particular removing $o(\lambda^2/\log\lambda)$
dimensions cannot leave a fixed positive or positive polynomial-in-
$\lambda^{-1}$ gap. The complete spectral projection onto
$[-2\eta_a,2\eta_a]$ has rank at least $2J+1$, with at least
$J+1$ even and $J$ odd eigenvalues counted with multiplicity.
\end{corollary}
\begin{proof}
Dimension gives the indicated unit vector, and the complete residual
gives its form bound. If the spectral projection had smaller rank,
a unit vector in the block orthogonal to its range would have
operator norm at least $2\eta_a$, contradicting
\eqref{ns46-a2-eq-large-residual}. Apply the same argument to each
parity block. The complete self-adjoint operators have compact
resolvent, as already established in the manuscript.
\end{proof}
With an actual even source constraint already imposed, fewer than
$J$ additional independent even constraints still leave such a vector;
fewer than $J$ odd constraints do so in the odd sector. More precisely
the relevant number is the rank of the constraint functionals restricted
to the indicated block. No identification with a sampler image or a
finite Fourier head is made.

\paragraph{What remains missing.}
The block construction supplies many complete near-zero directions and
their full residual; it supplies neither their energy signs nor a lower
bound on the infinite-dimensional orthogonal complement. The remaining
signed estimate is exactly
\[
 q_a[y]\ge-C_*\|y\|^2\quad
 (y\in\mathcal D(q_a)\cap\operatorname{Ran}(I-P_G)),
 \qquad C_*<\infty\ \hbox{uniformly on a cofinal family},
\]
in both parities. Together with the residual, the existing gap-free
transfer in v1.18/v1.36 would give a full uniform floor. Rewriting this
estimate in block coordinates is not new arithmetic input. Increasing
the block's rank does not bound negative eigenvalues remaining outside
it. The result is a stronger necessary rank obstruction for a positive
gap strategy; it does not exclude gap-free block methods, produce a
negative Weil direction, prove G2 or RH, or close the signed-floor route.
